\documentclass[11pt]{article}

\def\chapitre{20}
\def\pagetitle{Calcul différentiel.}
\input{setup.tex}

\begin{document}

\input{title.tex}

Dans tout le chapitre, $E$ et $F$ sont des $\R$-evdf. On rappelle quelques résultats utiles en \bf{dimension finie} :\\
$\hra$ Toutes les normes sont équivalentes.\\
$\hra$ $K$ est un compact ssi il est fermé et borné.\\
$\hra$ L'image continue d'un compact est compacte.\\
$\hra$ Une fonction continue sur un compact est bornée et atteint ses bornes.\\
$\hra$ Les applications multi-linéaires sont continues, donc aussi lipschitziennes.\\
$\hra$ Les applications coordonnées sont linéaires donc continues et lipschitziennes.\\
$\hra$ Tous les polynômes à plusieurs variables sont continus.

\end{document}
